\documentclass[11pt,a4paper]{article}

\usepackage[T1]{fontenc}
\usepackage[utf8]{inputenc}
\usepackage{lmodern}
\usepackage[scaled=0.95]{inconsolata}
\usepackage{microtype}
% Typewriter text: print "--" as two hyphens (command-line options), not as a dash.
\DisableLigatures[-]{encoding = *, family = tt*}
\usepackage[margin=2.4cm]{geometry}
\usepackage{xcolor}
\usepackage{listings}
\usepackage{booktabs}
\usepackage{tabularx}
\usepackage{enumitem}
\usepackage{caption}
\usepackage[most]{tcolorbox}
\usepackage{fancyhdr}
\usepackage[hidelinks]{hyperref}

% ---------------------------------------------------------------------------
% Text elements are black and white. Color is used only for the syntax
% highlighting of code, on a transparent background.
% ---------------------------------------------------------------------------
\definecolor{kw}{RGB}{0,70,160}       % keywords
\definecolor{cmt}{RGB}{90,110,90}     % comments
\definecolor{str}{RGB}{160,60,30}     % strings
\definecolor{lnum}{gray}{0.55}        % line numbers

% Notes and cautions: a black rule on the left, no fill.
\newtcolorbox{note}{enhanced,breakable,sharp corners,colback=white,colframe=black,
  boxrule=0pt,leftrule=0.8pt,left=8pt,right=0pt,top=2pt,bottom=2pt,
  before upper={\textbf{Note.}\ }}
\newtcolorbox{warning}{enhanced,breakable,sharp corners,colback=white,colframe=black,
  boxrule=0pt,leftrule=0.8pt,left=8pt,right=0pt,top=2pt,bottom=2pt,
  before upper={\textbf{Caution.}\ }}

% ---------------------------------------------------------------------------
% Listings: transparent background, thin black rules above and below
% ---------------------------------------------------------------------------
\lstset{
  basicstyle=\ttfamily\small,
  frame=lines, framerule=0.4pt, rulecolor=\color{black},
  breaklines=true, breakatwhitespace=false,
  postbreak=\mbox{$\hookrightarrow$\space},
  columns=fullflexible, keepspaces=true,
  showstringspaces=false,
  xleftmargin=0pt, xrightmargin=0pt,
  aboveskip=8pt, belowskip=8pt,
  captionpos=b,
  upquote=true,
}

% Shell sessions: lines that start with "$" are commands; the other lines are output.
\lstdefinestyle{shell}{
  language=bash,
  keywordstyle=\color{kw},
  commentstyle=\color{cmt}\itshape,
  stringstyle=\color{str},
  morekeywords={sudo,apt-get,opam,eval,jasminc,jasmin2ec,jasmin-ct,gcc,curl,unzip,df,du,npm,rm,make},
}
% Program output.
\lstdefinestyle{output}{
  language={}, basicstyle=\ttfamily\footnotesize,
}

\lstdefinelanguage{Jasmin}{
  morekeywords={fn,export,inline,return,reg,stack,inline,int,bool,u8,u16,u32,u64,
    u128,u256,for,to,downto,while,if,else,param,global,require,from},
  sensitive=true,
  morecomment=[l]{//},
  morecomment=[s]{/*}{*/},
  morecomment=[s]{\#[}{]},
  morestring=[b]",
}
\lstdefinestyle{jasmin}{
  language=Jasmin, keywordstyle=\color{kw}\bfseries,
  commentstyle=\color{cmt}\itshape, stringstyle=\color{str},
  numbers=left, numberstyle=\tiny\color{lnum}, numbersep=6pt,
}

\lstdefinelanguage{EasyCrypt}{
  morekeywords={require,import,from,lemma,proof,qed,module,proc,var,return,
    hoare,equiv,phoare,op,type,axiom,theory,end,clone,with,as,if,then,else,
    forall,exists,fun,by,smt,ring,rnd,skip,wp,rewrite,have,split,first,move,
    admit,res,true,false},
  sensitive=true,
  morecomment=[s]{(*}{*)},
  morestring=[b]",
}
\lstdefinestyle{easycrypt}{
  language=EasyCrypt, keywordstyle=\color{kw}\bfseries,
  commentstyle=\color{cmt}\itshape, stringstyle=\color{str},
  numbers=left, numberstyle=\tiny\color{lnum}, numbersep=6pt,
}
\lstdefinestyle{asm}{
  language=[x86masm]Assembler, morekeywords={adds,adcs,adc,ldr,str,movz,ret},
  keywordstyle=\color{kw}, commentstyle=\color{cmt},
}
\lstdefinestyle{c}{
  language=C, keywordstyle=\color{kw}\bfseries,
  commentstyle=\color{cmt}\itshape, stringstyle=\color{str},
  numbers=left, numberstyle=\tiny\color{lnum}, numbersep=6pt,
}

\newcommand{\cmd}[1]{\lstinline[style=shell,frame=none]{#1}}
\newcommand{\file}[1]{\texttt{#1}}

\pagestyle{fancy}
\fancyhf{}
\lhead{\small Jasmin and EasyCrypt on the Raspberry Pi 5}
\rhead{\small\thepage}

\title{\textbf{Installation and Testing of Jasmin and EasyCrypt\\on the Raspberry Pi 5}\\[4pt]
  \large A reproducible procedure for the formal verification of cryptographic code on AArch64}
\author{}
\date{October 21, 2026}

\begin{document}
\maketitle
\thispagestyle{empty}

\begin{abstract}
\noindent
This guide describes, step by step, how to install the \emph{Jasmin}
compiler (v2026.09.0) and the \emph{EasyCrypt} proof assistant (r2026.07)
on a Raspberry Pi~5 (8\,GB RAM, Debian~12 ``Bookworm'', AArch64), and how
to confirm with four tests, comprising 11 checks, that the toolchain
works. It is written for readers
who are familiar with the Linux command line but have not used Jasmin,
EasyCrypt, OCaml or Coq before; it explains each concept when it first
appears. Each executed command appears in a listing together with the
relevant part of its output. The tests demonstrate four capabilities of the
toolchain:
(i)~the compilation of a Jasmin program to native ARMv8-A assembly and its
execution against a C reference implementation;
(ii)~the verification of constant-time security with \texttt{jasmin-ct};
(iii)~the checking of EasyCrypt proofs, including a probabilistic
relational proof; and
(iv)~the translation of Jasmin code to EasyCrypt with \texttt{jasmin2ec}
and the proof of its functional correctness. All 11 checks pass.
The guide also documents four problems that arose during the installation
(a full SD card, a mislabeled Z3 release file, glibc incompatibilities, and
an operand restriction of the ARMv8-A back end) and their resolutions.
\end{abstract}

\tableofcontents
\newpage

%===========================================================================
\section{Introduction}
%===========================================================================

\subsection{Purpose and audience}

Cryptographic software must be correct for every input and must not reveal
secrets through side channels such as its execution time. Testing alone
cannot establish such properties, because it can examine only a small
fraction of the possible inputs. \emph{Formal verification} establishes them with machine-checked
mathematical proofs. Jasmin and EasyCrypt are two tools that make this
approach practical for cryptographic code. This guide installs both tools
on a Raspberry Pi~5 and demonstrates their use on small examples. It is the
first of three documents in this repository; the second document tests the
ARMv8-A back end of the Jasmin compiler, and the third applies both tools to
a post-quantum key encapsulation mechanism.

The guide assumes familiarity with the Linux command line (shell commands,
packages, environment variables). It does not assume prior knowledge of the
programming language OCaml, of the proof assistant Coq, or of Jasmin and
EasyCrypt. A glossary of the technical terms is provided in
Appendix~\ref{app:glossary}.

\subsection{The tools}

\textbf{Jasmin} is a programming language and a compiler for cryptographic
implementations that must be both fast and trustworthy (\emph{high
assurance}). Jasmin programs are written at the level of assembly: they name
registers, processor flags and memory accesses explicitly, so that the
programmer controls the generated machine code. The Jasmin compiler is
\emph{formally verified}: a proof, checked by the proof assistant Coq,
establishes that the assembly code that the compiler generates has the same
behavior as the source program. The Jasmin toolchain also includes checkers
for \emph{memory safety} (the absence of accesses outside the allocated
memory) and for
\emph{constant-time security}. A program is constant-time if secret data
influence neither its control flow (which branches it takes) nor the memory
addresses that it accesses; otherwise, an attacker who measures the
execution time can learn information about the secrets. \emph{Speculative}
constant-time security extends this property to the speculative execution
of modern processors (attacks of the Spectre family).

\textbf{EasyCrypt} is an interactive proof assistant for cryptographic
proofs: the user states a property, and EasyCrypt checks a proof of it
that the user writes step by step. EasyCrypt reasons about programs written
in a simple imperative language that may draw random values. It supports
\emph{Hoare logic}, which relates the states before and after the execution
of a program, and \emph{probabilistic relational Hoare logic} (pRHL),
which relates the output distributions of two programs; pRHL is the basis of
game-based security proofs, in which a real system is shown to be
indistinguishable from an idealized one. EasyCrypt delegates routine
arithmetic proof steps to automatic provers, called \emph{SMT solvers}
(Alt-Ergo, CVC5, Z3), through the intermediate platform \emph{Why3}.

The two tools are designed to work together. The translator
\texttt{jasmin2ec} converts a Jasmin program into an EasyCrypt module, on
which functional correctness and security properties can then be proved.
Because the compiler preserves the semantics of the program, these
properties also hold for the generated assembly code.

Jasmin v2026.09.0 introduced a \textbf{back end for ARMv8-A (AArch64)},
the 64-bit architecture of the processor of the Raspberry Pi~5 (a
Cortex-A76); a back end is the part of a compiler that generates code for
one particular architecture. Jasmin code can therefore be compiled,
executed, tested and verified on the same machine. This guide demonstrates each of these steps.

\subsection{Overview of the procedure}

Table~\ref{tab:roadmap} lists the steps of the installation. Most of the
time is spent compiling software from source; the commands themselves are
few, and Appendix~\ref{app:commands} collects them in one listing.

\begin{table}[h]
\centering\small
\begin{tabularx}{\textwidth}{@{}clXr@{}}
\toprule
\textbf{Step} & \textbf{Section} & \textbf{Action} & \textbf{Time} \\
\midrule
1 & \ref{sec:prep} & Free storage and check memory & a few min \\
2 & \ref{sec:apt} & Install system packages with \texttt{apt-get} & 2 min \\
3 & \ref{sec:opam} & Initialize opam and create an OCaml environment (switch) & 12 min \\
4 & \ref{sec:jasmin} & Install Jasmin with opam & 16 min \\
5 & \ref{sec:easycrypt} & Install EasyCrypt and the SMT solvers, and configure them & 57 min \\
6 & \ref{sec:tests} & Run the four tests (11 checks) & 15 s \\
\bottomrule
\end{tabularx}
\caption{Steps of the installation and their duration on the Raspberry Pi~5.}
\label{tab:roadmap}
\end{table}

\subsection{Conventions}

In shell listings, a line that begins with \texttt{\$} is a command that the
user types; the remaining lines are the output of the command, which is
sometimes abridged. The marker \texttt{...} indicates an omission. A
comment that begins with \texttt{\#} explains a command and is not part of
it. The example files are located in the directory \file{examples/} next to
this document; file paths in this guide are relative to the directory
\file{install-guide/} unless they begin with \texttt{/} or
\texttt{\textasciitilde}.

%===========================================================================
\section{Test platform}
%===========================================================================

Table~\ref{tab:platform} lists the hardware and software with which every
command of this guide was executed. Other versions may work, but they have
not been tested.

\begin{table}[h]
\centering
\begin{tabularx}{0.9\textwidth}{@{}lX@{}}
\toprule
\textbf{Item} & \textbf{Value} \\
\midrule
Board            & Raspberry Pi 5 Model B Rev 1.0, 4 cores (Cortex-A76), 7.9\,GiB RAM \\
Storage          & 64\,GB microSD card (\file{/dev/mmcblk0p2}, ext4) \\
Operating system & Debian GNU/Linux 12 (bookworm), Raspberry Pi OS kernel 6.12.93+rpt-rpi-2712, aarch64 \\
C compiler       & GCC 12.2.0 (Debian 12.2.0-14+deb12u1) \\
opam / OCaml     & opam 2.1.2 (Debian package), OCaml 4.14.2 (switch \texttt{jasmin-ec}) \\
Jasmin           & 2026.09.0 (\texttt{jasminc}, \texttt{jasmin2ec}, \texttt{jasmin-ct}, \texttt{jasmin2tex}) \\
EasyCrypt        & r2026.07 (git tag, pinned in opam) \\
Why3             & 1.8.2 \\
SMT solvers      & Alt-Ergo 2.6.0 (opam), CVC5 1.0.3 (apt), Z3 4.8.12 (apt) \\
\LaTeX           & pdfTeX 3.141592653-2.6-1.40.24 (TeX Live 2022/Debian) \\
\bottomrule
\end{tabularx}
\caption{Hardware and software of the test platform.}
\label{tab:platform}
\end{table}

%===========================================================================
\section{Preparation: storage and memory}
\label{sec:prep}
%===========================================================================

A build of EasyCrypt and Jasmin from source requires several gigabytes of
storage: the final installation occupies 4.0\,GB, and the build process
needs additional temporary space. The command \texttt{df} reports the free
storage, and the command \texttt{free} reports the free memory:

\begin{lstlisting}[style=shell]
$ df -h ~ | tail -1
/dev/mmcblk0p2   58G   55G  514M 100% /
$ free -h
               total        used        free      shared  buff/cache   available
Mem:           7.9Gi       5.2Gi       323Mi        14Mi       2.9Gi       2.7Gi
Swap:          2.5Gi       2.5Gi       288Ki
\end{lstlisting}

On the test platform, the SD card was full (100\,\% used). Storage was
reclaimed by deleting only caches and other files that their programs
regenerate on demand: the caches of npm, apt and pip, and superseded
versions of the VS Code Remote server, which VS Code installs on the
Raspberry Pi when a remote editing session connects to it. The server
version in use and the most recently used one were retained. Readers who do
not use these programs can skip the corresponding commands; the general
rule is to delete only files that are known to be regenerated.

\begin{lstlisting}[style=shell]
$ ps aux | grep -o 'servers/Stable-[0-9a-f]*' | sort | uniq -c   # which server is in use?
     30 servers/Stable-6d62b1d5dace93a33b3a6d3eca1317bd5e15bcb1
$ cd ~/.vscode-server/cli/servers && rm -rf Stable-<unused hashes...>
$ npm cache clean --force
$ sudo apt-get clean
$ rm -rf ~/.cache/pip
$ df -h / | tail -1
/dev/mmcblk0p2   58G   47G  8.2G  86% /
\end{lstlisting}

In the second command, \texttt{<unused hashes...>} stands for the
directory names of the server versions that the first command does not
report as in use.

\begin{note}
Reserve at least \textbf{6--8\,GB} of free storage. The available memory
(1.6--3\,GB, with other programs running) sufficed for opam's default
parallelism of 3 build jobs on 4 cores. If the operating system terminates a
build because memory is exhausted (a message from the \emph{out-of-memory
killer} in the output of \texttt{dmesg}), rerun \cmd{opam install} with the
option \cmd{-j 1}, which builds one package at a time.
\end{note}

%===========================================================================
\section{System packages}
\label{sec:apt}
%===========================================================================

Jasmin and EasyCrypt are installed with opam (Section~\ref{sec:opam}), but
they require several programs and libraries from the operating system.
Table~\ref{tab:packages} explains the role of each Debian package; the
following command installs them all:

\begin{lstlisting}[style=shell]
$ sudo apt-get update
$ sudo DEBIAN_FRONTEND=noninteractive apt-get install -y \
    opam build-essential m4 pkg-config unzip bubblewrap \
    libgmp-dev libpcre3-dev zlib1g-dev libmpfr-dev libppl-dev autoconf \
    z3 cvc5
$ opam --version
2.1.2
$ z3 --version
Z3 version 4.8.12 - 64 bit
$ cvc5 --version | head -1
This is cvc5 version 1.0.3
\end{lstlisting}

\begin{table}[h]
\centering\small
\begin{tabularx}{\textwidth}{@{}lX@{}}
\toprule
\textbf{Package} & \textbf{Role} \\
\midrule
\texttt{opam} & package manager for OCaml software, through which Jasmin and EasyCrypt are installed \\
\texttt{build-essential}, \texttt{m4}, \texttt{pkg-config}, \texttt{autoconf}, \texttt{unzip} & C compiler, linker and build tools needed to compile packages from source \\
\texttt{bubblewrap} & sandboxing tool that opam uses by default (disabled in Section~\ref{sec:opam}) \\
\texttt{libgmp-dev}, \texttt{libmpfr-dev}, \texttt{libppl-dev} & arbitrary-precision arithmetic and polyhedra libraries, required by Apron, the abstract-interpretation library of Jasmin's safety checker \\
\texttt{libpcre3-dev}, \texttt{zlib1g-dev} & regular-expression and compression libraries, installed as a precaution; none of the opam packages of this guide requires them, so they may be omitted \\
\texttt{z3}, \texttt{cvc5} & two SMT solvers that EasyCrypt can call \\
\bottomrule
\end{tabularx}
\caption{Debian packages and their roles.}
\label{tab:packages}
\end{table}

\begin{warning}
The EasyCrypt documentation recommends Z3 4.12.x. The official release file
\texttt{z3-4.12.4-arm64-glibc-2.35.zip}, however, contains an \emph{x86-64}
program instead of an arm64 program, so that its execution fails:
\begin{lstlisting}[style=output]
x86_64-binfmt-P: Could not open '/lib64/ld-linux-x86-64.so.2': No such file or directory
\end{lstlisting}
All later arm64 releases for Linux (4.15.x--5.3.x) require version 2.38 of
the C library glibc, whereas Bookworm provides version 2.36. This
installation therefore retains Debian's Z3 4.8.12 and uses Alt-Ergo 2.6.0,
the version that EasyCrypt recommends, as the principal prover. All tests in
this guide pass with this configuration. The opam package \texttt{z3} can be
built from source instead, but the build takes a long time on the
Raspberry Pi.
\end{warning}

%===========================================================================
\section{opam and the OCaml switch}
\label{sec:opam}
%===========================================================================

Jasmin and EasyCrypt are written in OCaml and are distributed through
\emph{opam}, the package manager of the OCaml ecosystem. opam installs
packages into a \emph{switch}: an independent directory under
\file{\textasciitilde/.opam} with its own OCaml compiler and its own set of
packages. Separate switches allow incompatible versions of software to
coexist; the third document of this repository, for example, uses a second
switch for a development version of Jasmin.

This installation uses a dedicated switch named \texttt{jasmin-ec} with
OCaml 4.14.2. The first command initializes opam: the option \cmd{--bare}
prevents opam from creating a default switch, and the option
\cmd{--disable-sandboxing} prevents failures of \texttt{bubblewrap} in
containers and remote sessions. The second command creates the switch and
compiles the OCaml compiler from source. The third command activates the
switch in the current shell, by setting environment variables such as
\texttt{PATH}.

\begin{lstlisting}[style=shell]
$ opam init -y --bare --disable-sandboxing
...
[WARNING] Shell not updated in non-interactive mode: use --shell-setup
$ opam switch create jasmin-ec ocaml-base-compiler.4.14.2 -y
...
-> installed ocaml-base-compiler.4.14.2
-> installed ocaml-config.2
-> installed ocaml.4.14.2
Done.
$ eval $(opam env --switch=jasmin-ec)
$ ocaml -version
The OCaml toplevel, version 4.14.2
\end{lstlisting}

The compilation of OCaml took about 12 minutes on the Raspberry Pi~5. The
activation applies only to the current shell. To activate the switch in
every new shell, add the \texttt{eval} line to \file{\textasciitilde/.zshrc}
or \file{\textasciitilde/.bashrc}, or run \cmd{opam init --shell-setup}
once. All commands in the remainder of this guide assume an active switch.

%===========================================================================
\section{Installation of Jasmin}
\label{sec:jasmin}
%===========================================================================

\subsection{Selection of versions}

The Jasmin toolchain translates programs into EasyCrypt code
(Section~\ref{sec:test4}), and each Jasmin release generates code for one specific EasyCrypt release;
the two versions must therefore be selected together. The following commands
retrieve the latest releases from GitHub and the compatibility statement in
Jasmin's change log (\file{CHANGELOG.md}):

\begin{lstlisting}[style=shell]
$ curl -s 'https://api.github.com/repos/jasmin-lang/jasmin/releases?per_page=1' | grep tag_name
    "tag_name": "v2026.09.0",
$ curl -s 'https://api.github.com/repos/EasyCrypt/easycrypt/releases?per_page=1' | grep tag_name
    "tag_name": "r2026.10",
$ curl -s https://raw.githubusercontent.com/jasmin-lang/jasmin/v2026.09.0/CHANGELOG.md \
    | grep -n 'targets version' | head -1
42:- Extraction as EasyCrypt code targets version 2026.07
\end{lstlisting}

The latest EasyCrypt release (r2026.10) is thus newer than the release that
Jasmin 2026.09.0 supports (r2026.07). Table~\ref{tab:versions} records the
selection.

\begin{table}[h]
\centering
\begin{tabular}{@{}llll@{}}
\toprule
\textbf{Tool} & \textbf{Latest release} & \textbf{Installed} & \textbf{Reason} \\
\midrule
Jasmin    & v2026.09.0 & 2026.09.0 & first release with the ARMv8-A target \\
EasyCrypt & r2026.10   & r2026.07  & version targeted by Jasmin 2026.09.0 \\
\bottomrule
\end{tabular}
\caption{Selected versions.}
\label{tab:versions}
\end{table}

\subsection{Installation with opam}

The first command lists the versions of Jasmin that opam offers; the second
installs version 2026.09.0 and measures the duration with \texttt{time}:

\begin{lstlisting}[style=shell]
$ opam show jasmin --field=all-versions
2022.04.0  2022.09.0 ... 2026.03.2  2026.09.0
$ time opam install -y jasmin.2026.09.0
...
-> installed apron.v0.9.15
-> installed jasmin.2026.09.0
Done.
<><> conf-gmp-paths.1 installed successfully <><><><><><><><><><><><><><><><><>
=> header file found in /usr/include/aarch64-linux-gnu
=> library file found in /usr/lib/aarch64-linux-gnu
opam install -y jasmin.2026.09.0  755.07s user 96.10s system 87% cpu 16:14.07 total
\end{lstlisting}

The Coq proof of the compiler is not checked during this installation: the
opam package contains the OCaml code that was extracted from the Coq
development, so the installation does not require Coq. The build took
16~minutes.

\subsection{Verification of the installation}

The package installs four programs, which Table~\ref{tab:jasmintools}
describes. The following commands confirm that they are available and that
the compiler supports the ARMv8-A target (\texttt{armv8a}):

\begin{lstlisting}[style=shell]
$ which jasminc jasmin2ec jasmin-ct jasmin2tex
/home/code-j/.opam/jasmin-ec/bin/jasminc
/home/code-j/.opam/jasmin-ec/bin/jasmin2ec
/home/code-j/.opam/jasmin-ec/bin/jasmin-ct
/home/code-j/.opam/jasmin-ec/bin/jasmin2tex
$ jasminc -version
Jasmin Compiler 2026.09.0
$ jasminc -help | grep -A1 -- '-arch'
  -arch {x86-64|arm-m4|armv8a|riscv}
                             Select target arch (default is x86-64)
\end{lstlisting}

\begin{table}[h]
\centering\small
\begin{tabularx}{\textwidth}{@{}lX@{}}
\toprule
\textbf{Program} & \textbf{Function} \\
\midrule
\texttt{jasminc} & compiler: translates a Jasmin program (\texttt{.jazz}) into assembly for the selected target \\
\texttt{jasmin-ct} & constant-time checker: verifies that secret data influence neither branches nor memory addresses \\
\texttt{jasmin2ec} & translator: converts a Jasmin program into an EasyCrypt module \\
\texttt{jasmin2tex} & pretty-printer: typesets a Jasmin program in \LaTeX{} (not used in this guide) \\
\bottomrule
\end{tabularx}
\caption{Programs of the Jasmin package.}
\label{tab:jasmintools}
\end{table}

The package also installs Jasmin's EasyCrypt support library, called
\emph{eclib} (\file{JModel\_armv8a.ec}, \file{JWord.ec}, \dots), in the
directory \file{\$OPAM\_SWITCH\_PREFIX/lib/easycrypt/jasmin}. This library
defines, in EasyCrypt, the machine words and the instructions that the
translated programs use. EasyCrypt locates it automatically
(Section~\ref{sec:ec-config}).

%===========================================================================
\section{Installation of EasyCrypt}
\label{sec:easycrypt}
%===========================================================================

\subsection{Installation with opam}

The default opam repository does not provide EasyCrypt at the required
version. The command \texttt{opam pin} therefore associates the package name
\texttt{easycrypt} with a specific source, the git tag \texttt{r2026.07} of
the EasyCrypt repository; the subsequent \texttt{opam install} builds the
package from this source, together with the SMT solver Alt-Ergo~2.6.0:

\begin{lstlisting}[style=shell]
$ opam pin add -yn easycrypt https://github.com/EasyCrypt/easycrypt.git#r2026.07
[easycrypt.~dev] synchronised (git+https://github.com/EasyCrypt/easycrypt.git#r2026.07)
easycrypt is now pinned to git+https://github.com/EasyCrypt/easycrypt.git#r2026.07 (version ~dev)
$ time opam install -y easycrypt alt-ergo.2.6.0
...
$ opam list --installed easycrypt why3 alt-ergo bitwuzla-cxx
# Name       # Installed # Synopsis
alt-ergo     2.6.0       The Alt-Ergo SMT prover
bitwuzla-cxx 0.9.1       SMT solver for AUFBVFP (C++ API)
easycrypt    ~dev        pinned to version ~dev at git+https://github.com/EasyCrypt/easycrypt.git#r2026.07
why3         1.8.2       Why3 environment for deductive program verification
\end{lstlisting}

This step took about \textbf{57 minutes}. Most of the time was spent on the
C++ build of the SMT solver Bitwuzla (\texttt{bitwuzla-cxx}), on Why3 and
Alt-Ergo, and on a rebuild of Jasmin: the dependencies of EasyCrypt require
version 1.3.0 of the OCaml library \texttt{cmdliner} instead of version
2.1.1, so opam recompiled every installed package that depends on
\texttt{cmdliner}, including Jasmin.

\begin{note}
The terminal session was disconnected during this step. The
\texttt{opam install} process continued in the background and completed
normally. In such a case, check with \cmd{pgrep -af opam} whether an opam
process is still running before starting another installation in the same
switch: opam locks the switch while it works, and a second
\texttt{opam install} waits until the lock is released.
\end{note}

\subsection{Configuration of the SMT solvers}
\label{sec:ec-config}

EasyCrypt calls the SMT solvers through Why3, which must know where each
solver is installed. The command \texttt{easycrypt why3config} searches for
the installed solvers and writes this information to a configuration file;
it must be rerun after every installation or update of a solver. The command
\texttt{easycrypt config} then displays the resulting configuration:

\begin{lstlisting}[style=shell]
$ easycrypt why3config
...
Found prover Z3 version 4.8.12, OK.
Found prover Z3 version 4.8.12 (alternative: noBV)
10 prover(s) added
Save config to /home/code-j/.config/easycrypt/why3.conf
$ easycrypt config
git-hash: r2026.07
load-path:
  Jasmin@/home/code-j/.opam/jasmin-ec/lib/easycrypt/jasmin
  <system>@/home/code-j/.opam/jasmin-ec/lib/easycrypt/theories
  ...
why3 configuration file
  /home/code-j/.config/easycrypt/why3.conf
...
known provers: Alt-Ergo@2.6.0, Alt-Ergo[BV]@2.6.0, Alt-Ergo[counterexamples]@2.6.0,
  CVC5@1.0.3, CVC5[counterexamples]@1.0.3, CVC5[strings]@1.0.3,
  CVC5[strings+counterexamples]@1.0.3, Z3@4.8.12, Z3[counterexamples]@4.8.12, Z3[noBV]@4.8.12
\end{lstlisting}

Two entries of this output are relevant. The entry \texttt{known provers}
lists the solvers that EasyCrypt can invoke; several solvers appear more than
once, in different configurations. The entry \texttt{Jasmin@\dots} of the
load path shows that the Jasmin eclib is registered under the logical name
\texttt{Jasmin}; the EasyCrypt statement
\lstinline[style=easycrypt,numbers=none,frame=none]|from Jasmin require import JModel_armv8a.|
therefore finds the library without additional options.

%===========================================================================
\section{Tests}
\label{sec:tests}
%===========================================================================

The tests confirm that each component of the toolchain works, and they
introduce the basic use of each tool. Table~\ref{tab:tests} summarizes what
each test demonstrates. The test files are located in \file{examples/}, and
a single script runs all tests (Section~\ref{sec:script}). The directory
has the following structure:

\begin{lstlisting}[style=output]
examples/
|-- run-tests.sh              # runs all 11 checks
|-- run-tests-output.txt      # output of the last run
|-- jasmin/
|   |-- add256.jazz           # 256-bit addition (ARMv8-A)
|   |-- test_add256.c         # C harness, 10^6 random tests
|   `-- ct_select.jazz        # constant-time and leaky selection
`-- easycrypt/
    |-- hello.ec              # SMT smoke test
    |-- otp.ec                # one-time pad, pRHL proof
    |-- wrong.ec              # negative test: must fail
    |-- CtSelect.ec           # generated by jasmin2ec
    `-- ct_select_proof.ec    # correctness proof of the Jasmin code
\end{lstlisting}

\begin{table}[h]
\centering\small
\begin{tabularx}{\textwidth}{@{}lXl@{}}
\toprule
\textbf{Test} & \textbf{What the test demonstrates} & \textbf{Tools} \\
\midrule
1 & Jasmin code compiles to ARMv8-A assembly that runs on the Raspberry Pi and computes correct results & \texttt{jasminc}, \texttt{gcc} \\
2 & The constant-time checker accepts code without secret-dependent branches and rejects code with such a branch & \texttt{jasmin-ct} \\
3 & EasyCrypt accepts correct proofs, including a probabilistic one, and rejects an attempted proof of a false statement & \texttt{easycrypt} \\
4 & A Jasmin function, translated to EasyCrypt, can be proved functionally correct & \texttt{jasmin2ec}, \texttt{easycrypt} \\
\bottomrule
\end{tabularx}
\caption{Overview of the tests. The test script (Section~\ref{sec:script}) divides these four tests into 11 checks.}
\label{tab:tests}
\end{table}

%---------------------------------------------------------------------------
\subsection{Test 1: compilation of Jasmin to ARMv8-A and native execution}
%---------------------------------------------------------------------------

The Jasmin program of the first test adds two 256-bit integers. A 64-bit processor cannot
add such numbers in one instruction, so each number is represented by four
64-bit words, called \emph{limbs}, and the limbs are added from the least to
the most significant one. Each addition sets the processor's \emph{carry
flag} when its result exceeds 64 bits, and the next addition adds this
carry. Multi-precision addition of this kind is a basic building block of
elliptic-curve and other public-key arithmetic. The program manipulates the
carry flag explicitly, through the Boolean variable \texttt{cf}. The
\texttt{for} loop assigns the values 1, 2 and 3 to the compile-time variable
\texttt{i} of type \texttt{inline int} (the upper bound 4 is excluded), and
the compiler unrolls the loop completely.

\lstinputlisting[style=jasmin,caption={\file{examples/jasmin/add256.jazz}}]{examples/jasmin/add256.jazz}

The compiler translates the program into an assembly file for the ARMv8-A
target:

\begin{lstlisting}[style=shell]
$ cd examples/jasmin
$ jasminc -arch armv8a -o add256.s add256.jazz
$ cat add256.s
\end{lstlisting}

\begin{lstlisting}[style=asm,caption={AArch64 assembly generated by \texttt{jasminc} (\file{add256.s})}]
	.global	add256
	.p2align	2
	.type	add256, %function
add256:
	str 	x30, [sp, #-16]!
	ldr 	x3, [x1]
	ldr 	x4, [x2]
	adds	x3, x3, x4
	str 	x3, [x0]
	ldr 	x3, [x1, #8]
	ldr 	x4, [x2, #8]
	adcs	x3, x3, x4
	str 	x3, [x0, #8]
	ldr 	x3, [x1, #16]
	ldr 	x4, [x2, #16]
	adcs	x3, x3, x4
	str 	x3, [x0, #16]
	ldr 	x3, [x1, #24]
	ldr 	x4, [x2, #24]
	adcs	x3, x3, x4
	str 	x3, [x0, #24]
	movz	x0, #0
	movz	x1, #0
	adc 	x0, x1, x0
	ldr 	x30, [sp], #16
	ret
	.ident	"Jasmin Compiler 2026.09.0"
	.section	".note.GNU-stack", "", %progbits
\end{lstlisting}

The generated code follows AAPCS64, the procedure call standard of AArch64,
which defines how a function receives its arguments and returns its result:
the three pointer arguments arrive in the registers
\texttt{x0}--\texttt{x2}, and the final carry is returned in \texttt{x0}.
A C program can therefore call the function directly. The carry chain
compiles to one \texttt{adds} instruction (addition that sets the carry
flag) followed by \texttt{adcs} instructions (addition with carry). A C test
program compares the result with a portable reference implementation, which
computes with the 128-bit C type \texttt{unsigned \_\_int128}, on $10^6$
random inputs. Every seventh input is constructed so that the carry
propagates through all four limbs:

\lstinputlisting[style=c,caption={\file{examples/jasmin/test\_add256.c}}]{examples/jasmin/test_add256.c}

The C compiler links the C test program with the generated assembly file:

\begin{lstlisting}[style=shell]
$ gcc -O2 -Wall -o test_add256 test_add256.c add256.s
$ ./test_add256
add256: 1000000 random tests passed
\end{lstlisting}

The same source file also compiles for the other targets of Jasmin, for
example x86-64. In this case \texttt{jasminc} acts as a cross-compiler: it
generates code for a processor other than the one on which it runs, and the
Raspberry Pi cannot execute the result:

\begin{lstlisting}[style=shell]
$ jasminc -arch x86-64 -o add256_x86.s add256.jazz
\end{lstlisting}

\begin{warning}
The first version of the program ended with the statement
\lstinline[style=jasmin,numbers=none,frame=none]|_, carry += 0 + cf;|,
which the compiler rejected:
\begin{lstlisting}[style=output]
"add256.jazz", line 21 (2-21):
compilation error in function add256:
asmgen: instruction ADC_64 is given incompatible args.
Allowed args are:
  [[reg]; [reg]; [reg]]
\end{lstlisting}
The AArch64 instruction \texttt{ADC} accepts only registers as operands, not
a constant (an \emph{immediate}). Jasmin reports this restriction instead of
inserting the necessary register moves implicitly, because each Jasmin
instruction corresponds to exactly one machine instruction. The correction
places the constant in a register
(\lstinline[style=jasmin,numbers=none,frame=none]|zero = 0;|) and uses the
statement \lstinline[style=jasmin,numbers=none,frame=none]|_, carry += zero + cf;|.
\end{warning}

%---------------------------------------------------------------------------
\subsection{Test 2: constant-time verification with \texttt{jasmin-ct}}
%---------------------------------------------------------------------------

A conditional selection returns one of two values depending on a condition
bit. If the condition bit is secret and the selection is implemented with a
branch, the execution time and the memory accesses can depend on the
secret. The checker \texttt{jasmin-ct} verifies that secret data influence
neither the control flow nor the memory addresses of a program. For this
purpose, the annotation \texttt{\#[ct = "..."]} declares the \emph{security
type} of each argument and of the result: \texttt{public} (known to an
attacker) or \texttt{secret}. The test file defines two functions that
compute the same result: a branch-free selection, which combines the two values with a
mask computed from the condition bit, and a leaky selection, which branches
on the secret bit:

\lstinputlisting[style=jasmin,caption={\file{examples/jasmin/ct\_select.jazz}}]{examples/jasmin/ct_select.jazz}

The option \texttt{--slice} restricts the check to one function. The first
command checks only \texttt{ct\_select}; the second checks the whole file:

\begin{lstlisting}[style=shell]
$ jasmin-ct --arch=armv8a --slice=ct_select ct_select.jazz; echo "exit=$?"
exit=0
$ jasmin-ct --arch=armv8a ct_select.jazz; echo "exit=$?"

"ct_select.jazz", line 22 (6-7):
constant type checker: c has type secret it needs to be public
exit=1
\end{lstlisting}

The checker accepts \texttt{ct\_select} (exit code 0) and rejects
\texttt{leaky\_select} (exit code 1); its error message identifies the
offending condition \texttt{c == 1} in line 22. The option
\cmd{--speculative} additionally verifies speculative constant time
(Spectre v1).

%---------------------------------------------------------------------------
\subsection{Test 3: EasyCrypt proofs}
%---------------------------------------------------------------------------

An EasyCrypt file states \emph{lemmas}, each followed by a proof between the
keywords \texttt{proof} and \texttt{qed}. A proof consists of
\emph{tactics}, commands that transform the statement to be proved (the
\emph{goal}) into simpler goals until none remains. EasyCrypt rejects the
file if any goal remains open. This test checks three files.

\paragraph{SMT smoke test.} The first file verifies that the standard
library loads, that the tactic \texttt{smt()} can delegate goals to the SMT
solvers through Why3, and that the tactic \texttt{ring}, which proves
algebraic identities, is available:

\lstinputlisting[style=easycrypt,caption={\file{examples/easycrypt/hello.ec}}]{examples/easycrypt/hello.ec}

\paragraph{Probabilistic relational proof.} The one-time pad encrypts a
message bit \(m\) as \(m \oplus k\) with a uniformly random key bit \(k\).
It is perfectly secret: the ciphertext has the same distribution as a fresh
uniform bit, whatever the message. The file states this property in pRHL as
an equivalence (\texttt{equiv}) between two procedures, the real encryption
\texttt{OTP.enc} and an ideal procedure \texttt{Ideal.enc} that returns a
random bit. The postcondition \texttt{=\{res\}} requires that the results of
the two procedures be equal when their random choices are suitably paired;
such a pairing, called a \emph{coupling}, implies that the two results have
the same distribution. The tactic \texttt{rnd} establishes the coupling: it
pairs each key \(k\) of the real procedure with the value \(m \oplus k\) of
the ideal procedure. Because this pairing is a bijection between the two
sets of random choices, it preserves probabilities, and the two results
coincide:

\lstinputlisting[style=easycrypt,caption={\file{examples/easycrypt/otp.ec}}]{examples/easycrypt/otp.ec}

\paragraph{Negative test.} A proof checker is useful only if it rejects
false statements. The following file states a false lemma ($x + 1 = x$),
which EasyCrypt must reject:

\lstinputlisting[style=easycrypt,caption={\file{examples/easycrypt/wrong.ec}}]{examples/easycrypt/wrong.ec}

The command \texttt{easycrypt compile} checks a file and reports success
with exit code 0:

\begin{lstlisting}[style=shell]
$ cd examples/easycrypt
$ easycrypt compile -no-eco hello.ec; echo "exit=$?"
exit=0
$ easycrypt compile -no-eco otp.ec; echo "exit=$?"
exit=0
$ easycrypt compile -no-eco wrong.ec; echo "exit=$?"
[critical] [wrong.ec: line 5 (7-16)] cannot prove goal (strict)
exit=1
\end{lstlisting}

The option \texttt{-no-eco} disables the cache of previously checked
theories (\texttt{.eco} files) and thereby forces a complete check. On the
Raspberry Pi~5, the checks took 1.2\,s (\file{hello.ec}) and 2.6\,s
(\file{otp.ec}).

\begin{note}
\textbf{Inspection of proof goals without an editor.} When a tactic fails,
\texttt{easycrypt compile} reports only the error. To display the goal that
remains at a given line, pass the file up to that line to the command-line
interface of EasyCrypt:
\begin{lstlisting}[style=shell]
$ sed -n '1,12p' ct_select_proof.ec | easycrypt cli -emacs
...
Current goal (remaining: 2)
------------------------------------------------------------------------
a0 `&` - one `|` b0 `&` invw (- one) = a0
\end{lstlisting}
For regular proof development, an interactive editor is more convenient:
Proof General (Emacs) or the EasyCrypt extension for VS Code.
\end{note}

%---------------------------------------------------------------------------
\subsection{Test 4: translation from Jasmin to EasyCrypt and correctness proof}
\label{sec:test4}
%---------------------------------------------------------------------------

This test combines the two tools. The translator \texttt{jasmin2ec}
converts the constant-time function \texttt{ct\_select} of Test~2 into an
EasyCrypt module (\texttt{M}) with one procedure per translated function. The option
\texttt{--arch=armv8a} selects the eclib model of the ARMv8-A instructions,
and the option \texttt{-f} selects the function to translate:

\begin{lstlisting}[style=shell]
$ cd examples/jasmin
$ jasmin2ec --arch=armv8a -f ct_select -o ../easycrypt/CtSelect.ec ct_select.jazz
\end{lstlisting}

\lstinputlisting[style=easycrypt,caption={\file{examples/easycrypt/CtSelect.ec} (generated)}]{examples/easycrypt/CtSelect.ec}

A proof in Hoare logic then establishes that the branch-free code computes a
selection. The statement \texttt{hoare [ M.ct\_select : P ==> Q ]} asserts
that every execution of the procedure that starts in a state satisfying the
precondition \texttt{P} ends in a state satisfying the postcondition
\texttt{Q}. Here, the precondition requires \(c \in \{0, 1\}\), and the
postcondition requires that the result \texttt{res} equals \(a\) if
\(c = 1\) and \(b\) if \(c = 0\). The proof treats the two cases separately
and relies on the bit-level properties of the mask \(0 - c\), which consists
of either all ones or all zeros, using lemmas of \file{JWord.ec}:

\lstinputlisting[style=easycrypt,caption={\file{examples/easycrypt/ct\_select\_proof.ec}}]{examples/easycrypt/ct_select_proof.ec}

\begin{lstlisting}[style=shell]
$ cd ../easycrypt
$ easycrypt compile -no-eco ct_select_proof.ec; echo "exit=$?"
exit=0
\end{lstlisting}

The proof is checked in 6.5\,s on the Raspberry Pi~5; most of this time is
spent loading the Jasmin eclib. Together with Test~2, this result
establishes that \texttt{ct\_select} is constant-time (\texttt{jasmin-ct})
\emph{and} functionally correct (EasyCrypt). Because the Jasmin compiler is
verified, both properties also hold for the generated AArch64 code.

%---------------------------------------------------------------------------
\subsection{Execution of all tests}
\label{sec:script}
%---------------------------------------------------------------------------

The script \file{examples/run-tests.sh} rebuilds and checks all examples. It
performs 11 checks, which divide the four tests above into their individual
steps. For the two negative checks (the leaky selection and the false
lemma), the script verifies that the tool fails as intended, with exit
code~1:

\lstinputlisting[style=shell,caption={\file{examples/run-tests.sh}}]{examples/run-tests.sh}

\begin{lstlisting}[style=shell]
$ cd examples
$ time ./run-tests.sh | tee run-tests-output.txt
\end{lstlisting}

\lstinputlisting[style=output,caption={Output of \file{run-tests.sh} on the Raspberry Pi 5}]{examples/run-tests-output.txt}

The complete test suite runs in about 15\,s. Each line marked
\texttt{[PASS]} corresponds to one check; a failed check is marked
\texttt{[FAIL]} and makes the script exit with a nonzero code.

%===========================================================================
\section{Summary}
%===========================================================================

Table~\ref{tab:summary} summarizes the time and storage that the
installation required.

\begin{table}[h]
\centering\small
\begin{tabularx}{\textwidth}{@{}lrX@{}}
\toprule
\textbf{Step} & \textbf{Wall-clock time} & \textbf{Remarks} \\
\midrule
apt packages                          & $\approx$ 2 min  & opam, build tools, Z3, CVC5 \\
\texttt{opam init} and OCaml 4.14.2 switch & $\approx$ 12 min & compiler built from source \\
Jasmin 2026.09.0                      & 16 min          & including Apron, PPL bindings, Menhir \\
EasyCrypt r2026.07 and Alt-Ergo 2.6.0 & $\approx$ 57 min & including Bitwuzla (C++), Why3, rebuild of Jasmin \\
Test suite (4 tests, 11 checks)       & 15 s            & all checks pass \\
\midrule
Storage used by \file{\textasciitilde/.opam} & 4.0 GB  & switch \texttt{jasmin-ec}: 3.9 GB \\
\bottomrule
\end{tabularx}
\caption{Time and storage requirements on the Raspberry Pi 5 (8\,GB RAM, microSD storage).}
\label{tab:summary}
\end{table}

\paragraph{Recommendations.}
\begin{itemize}[leftmargin=*]
  \item \textbf{Select matching versions.} Install the EasyCrypt release
        that Jasmin's \file{CHANGELOG.md} names (``Extraction as EasyCrypt
        code targets version \dots''), not the latest EasyCrypt release.
  \item \textbf{Reclaim storage before the build.} A margin of 6--8\,GB is
        sufficient. On a full disk, builds fail with misleading error
        messages.
  \item \textbf{Verify prover binaries for arm64.} The upstream Z3 builds
        for arm64 are either mislabeled (4.12.4) or require
        glibc~$\geq$~2.38. On Bookworm, use Debian's Z3 together with
        Alt-Ergo from opam.
  \item \textbf{Observe the ARMv8-A operand rules.} Instructions such as
        \texttt{ADC} accept only register operands; when a program violates
        such a rule, \texttt{jasminc} reports the permitted operand forms.
  \item \textbf{Run long builds in a persistent session.} The \texttt{opam}
        process continues after a disconnection, but a terminal multiplexer
        such as \texttt{tmux} or \texttt{screen} keeps its output visible
        and allows reconnecting to it.
\end{itemize}

The second document of this repository (\file{../armv8-tests}) builds on
this installation: it tests the ARMv8-A back end of Jasmin systematically on
the Raspberry Pi~5.

%===========================================================================
\appendix
\section{Complete command sequence}
\label{app:commands}
%===========================================================================

\begin{lstlisting}[style=shell]
# 1. System packages
$ sudo apt-get update
$ sudo apt-get install -y opam build-essential m4 pkg-config unzip bubblewrap \
    libgmp-dev libpcre3-dev zlib1g-dev libmpfr-dev libppl-dev autoconf z3 cvc5

# 2. opam and OCaml
$ opam init -y --bare --disable-sandboxing
$ opam switch create jasmin-ec ocaml-base-compiler.4.14.2 -y
$ eval $(opam env --switch=jasmin-ec)

# 3. Jasmin
$ opam install -y jasmin.2026.09.0

# 4. EasyCrypt and Alt-Ergo
$ opam pin add -yn easycrypt https://github.com/EasyCrypt/easycrypt.git#r2026.07
$ opam install -y easycrypt alt-ergo.2.6.0
$ easycrypt why3config
$ easycrypt config

# 5. Tests
$ cd install-guide/examples && ./run-tests.sh
\end{lstlisting}

%===========================================================================
\section{Glossary}
\label{app:glossary}
%===========================================================================

\begin{small}
\begin{description}[leftmargin=0pt,itemsep=2pt]
  \item[AAPCS64] Procedure Call Standard for the Arm 64-bit Architecture: the
        rules by which functions pass arguments, return results and preserve
        registers.
  \item[AArch64, ARMv8-A] The 64-bit execution state and the version of the
        Arm architecture that the Cortex-A76 of the Raspberry Pi~5
        implements; used interchangeably in this guide.
  \item[Back end] The part of a compiler that generates code for one
        particular processor architecture (target).
  \item[Carry flag] A processor flag that records whether the last addition
        produced a result too large for its register.
  \item[Constant time] The property that secret data influence neither the
        control flow nor the memory addresses of a program, so that its
        execution time reveals nothing about the secrets.
  \item[Coq] A proof assistant in which the correctness of the Jasmin
        compiler is proved.
  \item[Coupling] In pRHL: a pairing of the random choices of two programs
        that preserves their probabilities.
  \item[Cross-compiler] A compiler that generates code for a processor
        other than the one on which it runs.
  \item[eclib] The EasyCrypt library of Jasmin, which models machine words
        and instructions for programs translated by \texttt{jasmin2ec}.
  \item[Formal verification] The establishment of properties of a program by
        machine-checked mathematical proof, for all inputs.
  \item[Hoare logic] A logic for statements of the form ``if the precondition
        holds before the execution, the postcondition holds after it''.
  \item[Immediate] A constant operand that is encoded directly in an
        instruction.
  \item[Lemma, tactic, goal] In EasyCrypt: a stated property, a proof step,
        and a statement that remains to be proved.
  \item[Limb] One machine word of a multi-word integer.
  \item[opam, switch, pin] The OCaml package manager; an independent
        installation of an OCaml compiler with its packages; the association
        of a package name with a specific source.
  \item[pRHL] Probabilistic relational Hoare logic: a logic that relates the
        output distributions of two probabilistic programs.
  \item[Security type] In \texttt{jasmin-ct}: the annotation of a value as
        \texttt{public} or \texttt{secret}.
  \item[SMT solver] An automatic prover for logical formulas over arithmetic
        and other theories (here Alt-Ergo, CVC5 and Z3).
  \item[Why3] A platform through which EasyCrypt calls the SMT solvers.
\end{description}
\end{small}

%===========================================================================
\section{References}
%===========================================================================

\begin{itemize}[leftmargin=*]
  \item Jasmin: \url{https://github.com/jasmin-lang/jasmin}
  \item Jasmin documentation: \url{https://jasmin-lang.readthedocs.io}
  \item EasyCrypt: \url{https://github.com/EasyCrypt/easycrypt}
  \item Formosa-Crypto Zulip chat (user community of Jasmin and EasyCrypt):\\
        \url{https://formosa-crypto.zulipchat.com}
\end{itemize}

\end{document}
